\section{Studi Kasus: Optimasi Kode Bercabang \textit{If-Else} dan Gerbang Logika}

Untuk menutup selubung keanggunan Bab Ekuivalensi dan Tautologi ini, kita membongkar gerbang realita sejati ke panggung nyata mengapa kita mati-matian tersiksa menghafal puluhan manipulasi Distributif dan De Morgan \cite{rosen2019}. Aljabar Logika modern mutlak menyumbangkan dua kucuran pilar utama pada semesta teknologi modern: Menyederhanakan Baris percabangan program (Software), dan membabat biaya pencetakan gerbang Silikon Semikonduktor (Hardware).

\subsection{1. Pemangkasan Dead-Code di Belantara IF-ELSE}

Pernahkah Anda membaca kodingan kolega sistem senior yang berjubel kusut mengurai puluhan persyaratan Boolean gila di gerbang validasi server Web?

\textbf{Skenario Kode Asli Mentah:}
\begin{lstlisting}[style=pythonstyle]
def cek_validitas_pinjaman(a_kerja, b_jaminan, c_bebas_utang):
  
  # Validasi kualifikasi nasabah (versi ruwet):
  if not (not a_kerja and b_jaminan):
      # Cek tambahan di dalam:
      if a_kerja and (b_jaminan or not c_bebas_utang):
          return True
  return False
\end{lstlisting}

Mari kita bedah arsitekturnya ke ranah Logika Proposisional murni:
Sang programmer malang mendikte filter dengan pola: $\neg(\neg A \wedge B) \wedge (A \wedge (B \vee \neg C))$

\textbf{Proses Operasi Bedah Forensik Penyederhanaan:}
\begin{align*}
    1.\ & \neg(\neg A \wedge B) \wedge A \wedge (B \vee \neg C) & \text{\parbox[t]{5.8cm}{(Bentuk utuh gerbang bercabang program)}} \\
    2.\ & (A \vee \neg B) \wedge A \wedge (B \vee \neg C) & \text{\parbox[t]{5.8cm}{(De Morgan pada ruas awal)}} \\
    3.\ & A \wedge (A \vee \neg B) \wedge (B \vee \neg C) & \text{\parbox[t]{5.8cm}{(Hukum Komutatif)}} \\
    4.\ & \mathbf{A} \wedge (B \vee \neg C) & \text{\parbox[t]{5.8cm}{(Absorpsi: $A \wedge (A \vee \neg B)$ terisap ke $A$)}}
\end{align*}

\textbf{Transformasi Kode Bersih Pasca-Penyulingan (Refactoring):}
\begin{lstlisting}[style=pythonstyle]
def cek_validitas_pinjaman(a_kerja, b_jaminan, c_bebas_utang):
  
  # Filter Ekuivalensi Ceria nan Minimalis:
  return a_kerja and (b_jaminan or not c_bebas_utang)
\end{lstlisting}
Bayangkan efisiensi putaran memori *Central Processing Unit (CPU)* yang mampu diselamatkan. Itulah harga jutaan Dollar di balik sihir Ekuivalensi proporsi diskrit mahasiswa.

\subsection{2. Aplikasi Memutilasi Rangkaian Gerbang Logika Digital (Integrated Circuits)}

Sejatinya logika proposisional adalah ibu kandung pemodelan perangkat keras Gerbang Sirkuit Digital (Hardware \textit{AND/OR/NOT Gates}). Misalkan Divisi Pabrik ditawari kontrak mematenkan *blueprint* cetak IC pemroses dengan skematik sinyal Input A, B mengalir menunggangi formula: 
$$ Q_{output} = (A \land B) \lor (A \land \neg B) $$

Manajer bodoh akan kalap memborong 2 bilah Chip Gerbang AND beralih 1 Gerbang NOT, menguras saldo anggaran menelan jejak kelistrikan panjang memampatkan ruang papan sirkuit mahal berlapis tembaga.

\textbf{Analisis Ilmuwan Pintar (Penyederhanaan Silang):}
\begin{align*}
    (A \wedge B) \vee (A \wedge \neg B) &\equiv A \wedge (B \vee \neg B) \quad \text{(Distributif)} \\
    &\equiv A \wedge \mathbf{T} \quad \text{(Hukum Negasi)} \\
    &\equiv \mathbf{A} \quad \text{(Identitas)}
\end{align*}

Hasil Ekuivalensi mengaum mengerikan: $Q_{output} \equiv A$.
Artinya, \textbf{BUANG} rencana gila memborong 3 cetakan gerbang silikon mahal! Sambungkan saja Tembaga Kabel Input A polosan menerabas memotong jalan ke pintu gerbang pelabuhan akhir Output! Total biaya perakitan menukik tajam menjebol Rp0 perak. Pabrik menang tender \cite{wiki_proplogic}.
